\section{Results}
\label{sec:results}

We evaluate full consensus ($\text{CSS}=1.0$) across six configurations spanning two model architectures and four tasks.
Table~\ref{tab:main_results} summarizes the outcomes.
$\text{CSS}=1.0$ fails to outperform the best individual circuit discovery method in every configuration tested.
The failure modes differ across settings: a circuit-size confound on IOI, trivial collapse to a single method on Arithmetic, a statistically non-significant margin on SVA, significant degradation on Factual Recall, and empty consensus on both Qwen-2.5 configurations.
We present each result below, ordered by increasing severity of failure.

\begin{table*}[t]
\centering
\small
\caption{$\text{CSS}=1.0$ evaluation across all configurations. CPR measured under mean ablation ($\uparrow$). $\dagger$\,Size-controlled at $N\!=\!16$ edges. $\ddagger$\,Bootstrap CI ($n\!=\!1000$). Best CPR per row in \textbf{bold}.}
\label{tab:main_results}
\begin{tabular}{@{}lrcccrccl@{}}
\toprule
Configuration & $|C_{\tau=1}|$ & CPR$_{\text{CSS}}$ & CPR$_{\text{best}}$ & Best & $\Delta$ & 95\% CI & $p$ & Outcome \\
\midrule
GPT-2 IOI$^\dagger$ & 16 & 0.177 & \textbf{0.183} & EP & $-$0.006 & [$-$0.032, 0.014] & .279 & No advantage \\
GPT-2 Arith. & 17 & \textbf{0.985} & \textbf{0.985} & EAP-IG & 0.000 & & & $\equiv$ ACDC \\
GPT-2 SVA$^\ddagger$ & 44 & \textbf{1.187} & 1.175 & ACDC & $+$0.013 & [0.005, 0.026] & .099 & Not significant \\
GPT-2 FR$^\ddagger$ & & & & & neg. & & .000 & Sig.\ worse \\
Qwen-IFR & 0 & $\emptyset$ & & & & & & Empty \\
Qwen-ACDC & 0 & $\emptyset$ & & & & & & Empty \\
\bottomrule
\end{tabular}
\end{table*}

\paragraph{GPT-2 IOI.}
The indirect object identification task~\cite{wang2023interpretability} on GPT-2 Small initially appears to favor cross-method consensus.
Under zero ablation without size control, $\text{CSS}=1.0$ achieves a CPR of 0.1095, well above ACDC's 0.0286~\cite{conmy2023towards}.
This comparison is misleading because it confounds circuit quality with circuit size.
$\text{CSS}=1.0$ retains only 7 edges (the intersection of ACDC's 14 \texttt{isinf}-binarized edges with EAP-IG and EP top-100 circuits~\cite{syed2024attribution}), while individual methods retain between 14 and 100 edges.
Smaller circuits score higher on CPR when they happen to contain high-importance edges, even if they miss critical circuit components.
The pairwise Jaccard similarities between methods are low (EAP-IG${\times}$EP${=}$0.31, EAP-IG${\times}$ACDC${=}$0.10, EP${\times}$ACDC${=}$0.09), confirming that the three methods identify substantially different edge sets from the same model.

To isolate the size confound, we conduct a controlled experiment where all circuits are evaluated at exactly $N{=}16$ edges (Figure~\ref{fig:size_controlled}).
At equal size, $\text{CSS}=1.0$ achieves a mean-ablation CPR of 0.177, identical to ACDC top-16 (0.177) and below both EAP-IG top-16 (0.182) and EP top-16 (0.183).
Bootstrap confidence intervals ($n{=}1000$) confirm the absence of any significant difference: the comparison with EP yields $\Delta = -0.006$ with 95\% CI $[-0.032, 0.014]$ ($p = 0.279$); the comparison with EAP-IG yields $\Delta = -0.005$ with CI $[-0.039, 0.023]$ ($p = 0.359$).
Under zero ablation, the result is starker: $\text{CSS}=1.0$ and ACDC top-16 produce identical CPR values of 0.0295, while EAP-IG and EP top-16 both score $-0.0585$.
The apparent advantage of consensus was entirely a circuit-size artifact.

\begin{figure}[t]
\centering
\includegraphics[width=\columnwidth]{fig2_size_controlled.pdf}
\caption{IOI size-controlled comparison ($N{=}16$ edges). $\text{CSS}=1.0$ and ACDC top-16 produce identical CPR under zero ablation. Under mean ablation, $\text{CSS}=1.0$ performs below EAP-IG and EP top-16, with no statistically significant differences.}
\label{fig:size_controlled}
\end{figure}

\paragraph{GPT-2 Arithmetic.}
On the arithmetic task~\cite{hanna2023how}, $\text{CSS}=1.0$ produces exactly 17 edges, identical to the ACDC circuit.
This collapse occurs because ACDC's \texttt{isinf} binarization retains only 17 edges, all of which rank within the top 100 for both EAP-IG and EP.
The consensus set $C_{\tau=1}$ is therefore equal to the ACDC set.
Mean-ablation CPR for $\text{CSS}=1.0$ is 0.985, matching EAP-IG (0.985) and exceeding EP (0.962).
The pairwise Jaccard between EAP-IG and EP is 0.695, but both overlap with ACDC at only 0.17, reflecting the extreme size asymmetry between ACDC's 17-edge circuit and the 100-edge circuits of the gradient-based methods.
This pattern reveals a structural flaw: when one method applies a strict binarization threshold that produces a much smaller circuit, the intersection collapses to that method's output.
Consensus did not combine independent evidence from three methods; it selected the smallest circuit and discarded the unique contributions of the other two.

\paragraph{GPT-2 SVA.}
Subject-verb agreement~\cite{linzen2016assessing} provides the most favorable case for $\text{CSS}=1.0$ in our study.
All three methods use top-100 binarization, avoiding the \texttt{isinf} asymmetry that drives collapse on other tasks, and the CSS distribution is relatively smooth: 60 edges at threshold 1/3, 54 at 2/3, and 44 at 1.0.
$\text{CSS}=1.0$ achieves a mean-ablation CPR of 1.187, nominally above the best single method (ACDC at 1.175, with EAP-IG and EP both at 1.170).
However, bootstrap testing ($n{=}1000$) yields $p = 0.099$, failing to reach significance at $\alpha = 0.05$; the 95\% CI for $\Delta$ is $[0.005, 0.026]$.
The pairwise Jaccard values (EAP-IG${\times}$EP${=}$0.653, ACDC${\times}$EAP-IG${=}$0.316, ACDC${\times}$EP${=}$0.418) indicate moderate but incomplete overlap between methods.
Even in this best case, a non-significant margin of $+$0.013 CPR does not justify the computational cost of running three independent circuit discovery methods instead of one.

\paragraph{GPT-2 Factual Recall.}
Factual recall presents a qualitatively different failure mode.
EAP-IG and EP produce nearly identical circuits with a pairwise Jaccard of 0.82, far above any other method pair in our experiments.
When two of three methods share 82\% of their edges, $\text{CSS}=1.0$ converges toward their shared set and marginalizes the third method (ACDC).
This resembles majority voting among correlated classifiers: high inter-method agreement eliminates the diversity that consensus is supposed to exploit.
Bootstrap testing confirms that $\text{CSS}=1.0$ performs significantly worse than the best individual method ($p = 0.000$).
Cross-method consensus does not merely fail to improve circuit quality on this task; it actively degrades performance by amplifying the shared biases of two highly correlated gradient-based methods at the expense of the structurally distinct ACDC circuit.

\paragraph{Qwen-2.5 IOI.}
Scaling to Qwen-2.5 0.5B~\cite{qwen2025qwen} with 179,388 total edges exposes the most extreme failure mode.
In Config~5 (EAP-IG, EP, IFR), the CSS distribution concentrates at low consensus levels: 200 edges at threshold 1/3, 100 at 2/3, and exactly zero at $\text{CSS}=1.0$.
The consensus circuit is empty.
Pairwise Jaccard values explain this outcome: EAP-IG${\times}$EP${=}$0.54, but EAP-IG${\times}$IFR${=}$0.00 and EP${\times}$IFR${=}$0.00.
IFR identifies a completely disjoint edge set from the gradient-based methods, so unanimous agreement across all three is impossible regardless of threshold.

Replacing IFR with ACDC (Config~6) does not resolve the problem.
$\text{CSS}=1.0$ again yields zero edges, with maximum pairwise Jaccard at just 0.043 and $\text{CSS}{\geq}2/3$ producing only 66 edges.
At this model scale, methods disagree so thoroughly that no single edge receives support from all three.
The $\tau{=}1.0$ criterion produces a vacuous result, and even relaxing the threshold to $\tau{=}2/3$ recovers only a tiny fraction of the edge space.
These empty-consensus results demonstrate that cross-method disagreement worsens with model scale, as larger computation graphs provide more room for methods to diverge.

\paragraph{Within-method consensus (CIRCUS).}
One might expect that cross-method consensus fails because it compares fundamentally different algorithmic families.
CIRCUS~\cite{parekh2026circus} tests this alternative by replacing method diversity with threshold diversity: it runs a single method at multiple binarization thresholds and intersects the resulting circuits.
On IOI, CIRCUS-core produces 15 edges with a zero-ablation CPR of $-0.108$ and a CMD of 1.134, indicating substantial divergence from the full model's behavior.
On SVA, CIRCUS-core achieves a mean-ablation CPR of 0.994, comparable to individual methods at similar circuit sizes but not superior to any of them.
In neither case does within-method consensus outperform the constituent method run at a single well-chosen threshold.
The core issue is that threshold stability and functional importance are distinct properties of an edge.
An edge that appears consistently across binarization thresholds may simply occupy a stable position in the method's ranking, not a functionally critical role in the circuit.
Consensus, whether computed across methods or across thresholds, conflates these two properties and provides no guarantee of recovering edges that are genuinely important to the model's computation.
